\chapter{Lesson 10: School Life}

\section{Vocabulary}

\begin{center}
\begin{tabular}{llll}
\toprule
Characters & Pinyin & Part of speech & Meaning \\
\midrule
\zh{日记} & rìjì & noun & diary \\
\zh{起床} & qǐ chuáng & verb-object & to get up \\
\zh{洗澡} & xǐ zǎo & verb-object & to take a shower/bath \\
\zh{早饭} & zǎofàn & noun & breakfast \\
\zh{午饭} & wǔfàn & noun & lunch \\
\zh{晚饭} & wǎnfàn & noun & dinner \\
\zh{教室} & jiàoshì & noun & classroom \\
\zh{宿舍} & sùshè & noun & dormitory \\
\zh{食堂} & shítáng & noun & cafeteria \\
\zh{先} & xiān & adverb & first \\
\zh{再} & zài & adverb & then, again (future) \\
\zh{然后} & ránhòu & conjunction & afterwards \\
\zh{最后} & zuìhòu & adverb & finally \\
\zh{就} & jiù & adverb & as early as, right away \\
\zh{跟} & gēn & preposition & with \\
\zh{一起} & yìqǐ & adverb & together \\
\zh{每天} & měi tiān & & every day \\
\zh{电脑} & diànnǎo & noun & computer \\
\zh{上网} & shàng wǎng & verb-object & to go online \\
\zh{同学} & tóngxué & noun & classmate \\
\zh{专业} & zhuānyè & noun & major \\
\zh{学期} & xuéqī & noun & semester \\
\zh{最近} & zuìjìn & adverb & recently \\
\zh{能} & néng & modal & can, be able to \\
\zh{会} & huì & modal & can (learned skill); will \\
\zh{可以} & kěyǐ & modal & may, can (permitted) \\
\zh{希望} & xīwàng & verb & to hope \\
\zh{信} & xìn & noun & letter \\
\bottomrule
\end{tabular}
\end{center}

\section{The full Chinese word order}

All the pieces from earlier lessons now fit into one template.

\begin{center}
\begin{tikzpicture}[slot/.style={draw, thick, rounded corners, minimum height=1.6cm, minimum width=1.9cm, align=center, font=\small}]
  \node[slot, fill=blue!10] (s) at (0,0) {Subject\\[2pt] \zh{我}};
  \node[slot, fill=purple!10] (t) at (2.5,0) {Time\\[2pt] \zh{每天早上}\\ \zh{八点}};
  \node[slot, fill=gray!15] (a) at (5.0,0) {Adverb\\[2pt] \zh{常常}};
  \node[slot, fill=green!12] (p) at (7.5,0) {Place / \zh{跟}\\[2pt] \zh{跟同学}\\ \zh{在食堂}};
  \node[slot, fill=orange!18] (v) at (10.0,0) {Verb\\[2pt] \zh{吃}};
  \node[slot, fill=yellow!25] (o) at (12.5,0) {Object\\[2pt] \zh{早饭}};
  \foreach \a/\b in {s/t,t/a,a/p,p/v,v/o} \draw[->, thick] (\a) -- (\b);
  \node[font=\small, align=center] at (6.25,-1.35) {``I often have breakfast with classmates in the cafeteria at 8 every morning.''};
\end{tikzpicture}
\end{center}

\begin{important}[Time -- (manner/companion) -- place -- action]
Subject, then \emph{when}, then adverbs, then \emph{with whom} and \emph{where},
then the verb and its object. English is almost the mirror image. The time phrase may
also move in front of the subject (\zh{每天早上八点，我\ldots}), but place, companion
and adverbs stay between subject and verb.
\end{important}

\section{Grammar}

\subsection{Sequencing: \zh{先\ldots{}再\ldots{}然后\ldots{}最后}}

\begin{itemize}
  \item \zh{我先洗澡，再吃早饭。} I shower first, then eat breakfast.
  \item \zh{我们先去图书馆，然后去食堂吃午饭，最后回宿舍。}
\end{itemize}
\zh{再} here is ``then'' for actions not yet done (plans and habits); for past sequences
use \zh{然后} or \zh{就}.

\subsection{\zh{跟\ldots{}一起}}

\zh{跟} + person (+ \zh{一起}) + verb: ``do something with someone''.
\begin{itemize}
  \item \zh{我常常跟室友一起去打球。} I often play ball with my roommate.
  \item \zh{你想跟我们一起去吗？}
\end{itemize}

\subsection{\zh{就} vs \zh{才}}

Both come right before the verb, and they express the speaker's attitude to timing.

\begin{center}
\begin{tabular}{lll}
\toprule
 & \zh{就} jiù & \zh{才} cái \\
\midrule
Meaning & earlier / sooner than expected & later than expected \\
\zh{了}? & usually with \zh{了} & never with \zh{了} \\
Example & \zh{他六点就起床了。} & \zh{他十点才起床。} \\
 & He was up at six already. & He didn't get up until ten. \\
\bottomrule
\end{tabular}
\end{center}

\begin{center}
\begin{tikzpicture}[x=1.1cm]
  \draw[thick, -{Stealth}] (5,0) -- (11.5,0) node[right, font=\small] {time};
  \foreach \h in {6,7,...,11} {
    \draw (\h,0.1) -- (\h,-0.1);
    \node[below, font=\scriptsize] at (\h,-0.1) {\h:00};
  }
  \fill[gray] (8,0) circle (0.09);
  \node[above, font=\scriptsize, gray] at (8,0.15) {expected};
  \fill[green!50!black] (6,0) circle (0.11);
  \node[font=\small, green!40!black] at (6,-0.85) {\zh{就}};
  \draw[->, green!50!black, thick] (7.85,0.6) to[bend right=25] (6.15,0.6);
  \fill[red!70!black] (10,0) circle (0.11);
  \node[font=\small, red!70!black] at (10,-0.85) {\zh{才}};
  \draw[->, red!70!black, thick] (8.15,0.6) to[bend left=25] (9.85,0.6);
\end{tikzpicture}
\end{center}

\subsection{\zh{能}, \zh{会}, \zh{可以}}

All three translate as ``can''. A first look (revisited in CHINA 102R):
\begin{itemize}
  \item \zh{会}: a learned skill. \zh{我会说中文，也会开车。}
  \item \zh{能}: physical ability or circumstances. \zh{我明天有考试，今天晚上不能去。}
  \item \zh{可以}: permission. \zh{老师，我可以进来吗？}---\zh{可以。}
\end{itemize}
The negative of permission is usually \zh{不能} (``not allowed''), not \zh{不可以}
alone, though both are heard.

\subsection{Reporting a whole day: \zh{的时候} and \zh{以后} in narration}

\begin{itemize}
  \item \zh{上课的时候，老师说中文，我们也说中文。}
  \item \zh{吃了晚饭以后，我在宿舍做作业。}
\end{itemize}
Note \zh{V了 O 以后}: the \zh{了} marks the first action as complete before the second.

\section{A diary entry}

\begin{example}[\zh{我的一天}]
\zh{十月二十号 \ 星期三}

\zh{今天我七点就起床了。我先洗澡，再去食堂吃早饭。八点半我去教室上中文课。上课的时候，老师教了我们第十课的语法。第十课的语法不太难，可是生词有点儿多。中午我跟同学一起吃午饭，然后去图书馆用电脑做作业。下午我有数学课，五点才下课。吃了晚饭以后，我跟室友聊天、听音乐。我十一点半睡觉。}
\end{example}

\section{Writing a letter in Chinese}

\begin{aside}
A traditional Chinese letter (\zh{信}) has a fixed frame. Start with the addressee on
the first line followed by a colon: \zh{小王：} or \zh{亲爱的妈妈：} ``Dear Mom:''; then
\zh{你好！} on its own line. The body is indented two characters. Close with a wish such
as \zh{祝好！} ``All the best!'' or the formal two-line \zh{此致} / \zh{敬礼}, then
sign on the right with the date below. In traditional characters the closing is the
same, \zht{祝好}, while ``Dear'' \zh{亲爱的} is written \zht{親愛的}.
\end{aside}

\begin{example}[Short letter]
\zh{亲爱的小林：}

\zh{你好！好久不见，你最近怎么样？我这个学期很忙，因为我有五门课，还要学中文。中文有点儿难，可是很有意思。我希望明年能去北京学习。你什么时候有空？我们一起去喝咖啡吧。}

\zh{祝好！}

\hfill\zh{明}

\hfill\zh{十一月三日}
\end{example}

\section*{Review}
\begin{itemize}
  \item Word order: Subject -- Time -- Adverb -- \zh{跟}/Place -- Verb -- Object.
  \item \zh{先\ldots{}再\ldots{}然后\ldots{}最后} to sequence actions.
  \item \zh{就} (early, with \zh{了}) vs \zh{才} (late, no \zh{了}).
  \item \zh{会} skill, \zh{能} ability/circumstances, \zh{可以} permission.
  \item V \zh{了} O \zh{以后} for ``after doing''.
  \item Letter format: \zh{X：}, \zh{你好！}, body, \zh{祝好！}, signature and date.
\end{itemize}
